Chiplet-based systems are a modular design approach in which a complex system is divided into multiple smaller dies, called chiplets, that work together to implement the complete system functionality. Compared with traditional monolithic SoCs, this design style provides better flexibility, scalability, and hardware reuse, since functions such as computation, memory, and I/O can be implemented in separate chiplets and integrated according to application requirements. However, this modularity also makes communication a key challenge. Unlike monolithic chips, where data transfer relies on on-chip interconnects with relatively low latency and high bandwidth, chiplet-based systems must handle data movement across chiplet boundaries, which introduces additional communication overhead. As a result, the overall system performance depends strongly on how efficiently data can be transferred between chiplets \cite{11265742}.

UCIe has been developed as a standardized interface for chiplet communication. It addresses the physical and link-level transmission problem, that is, how data is transferred across die boundaries. However, for systems composed of multiple chiplets, physical connectivity alone is not sufficient. A higher-level communication mechanism is still required to determine how packets move from one chiplet to another in a correct and efficient way.

As the number of chiplets increases, communication in chiplet-based systems gradually evolves from simple point-to-point transfer into a networked communication problem. Recent studies show that chiplet-based systems increasingly adopt hierarchical communication structures, where each chiplet may contain its own local network, while an additional inter-chiplet network is used to connect different chiplets together.In this case, routing becomes a key issue because packets must traverse not only local paths inside a chiplet but also global paths across chiplet boundaries.

Routing in chiplet-based systems is more difficult than routing in conventional monolithic NoCs. In a monolithic NoC, routing algorithms are usually designed with complete knowledge of the whole topology. In chiplet-based systems, however, different chiplets may be designed independently and later integrated into one system. This means that global routing information may not be fully available during the design of each individual chiplet. As a result, even if every chiplet is deadlock-free on its own, the integrated system may still suffer from deadlocks caused by cyclic dependencies across chiplet boundaries \cite{chen2022design}.

For this thesis, the focus is not on the physical implementation details of chiplet interfaces, but on how data is routed efficiently across chiplets. Therefore, the key research issue is to design a lightweight and effective routing-oriented communication mechanism that supports correct and efficient data movement in a chiplet-based system.